\chapter{Radar and canopy height: what transfers}\label{ch:mangroves}
Mapping forest height is one of the commonest uses of radar. Here we test two kinds of radar measurement, backscatter strength and interferometric surface height, against GEDI lidar over the mangroves of the Indus Delta, and see what survives a fair test. Notebook: \texttt{mangroves.ipynb}.

\section{The question and the data}
\input{generated/nb_mangroves_1}
\section{What GEDI sees over short mangroves}
\input{generated/nb_mangroves_2}
\section{Each radar measurement against canopy height}
\input{generated/nb_mangroves_3}
\section{A fair test: spatial blocks}
\input{generated/nb_mangroves_4}
\section{Why VV falls as the canopy grows}
\input{generated/nb_mangroves_5}
\section{What you have learned}
\input{generated/nb_mangroves_6}
